\section{纯组合模型与巡检联合模型的联系}
\label{sec:01_03}

\subsection{从资源占用到执行见证}

纯组合模型回答的是：在给定候选任务及其资源需求后，哪些任务能够共同被接受。巡检联合模型还须回答：由哪些设备执行，何时出发，沿什么路径运动，在哪里完成观测，以及如何在满足能量和会合条件后结束任务。两类模型都包含选择决策，但对一个选择是否可行，所要求的证明不同。

沿用上一节的资源集合$V$、候选集合$E$和关联矩阵$A$，令$n=|E|$、$m=|V|$。选择变量$x_e\in\{0,1\}$表示是否接受候选$e\in E$，容量$b_v\in\mathbb Z_{\geq0}$表示资源$v\in V$可提供的占用份数。组合可行域为
\begin{equation}
 \mathcal F_{\mathrm c}
 =\left\{x\in\{0,1\}^{n}:Ax\leq b\right\}.
 \label{s0103:combinatorial_domain}
\end{equation}
其中，$A_{ve}\in\{0,1\}$记录资源与候选的关联，容量和占用份数均无量纲。这一表达以候选内容和资源含义已经固定为前提。例如，将某架无人机作为容量一的资源，表示所有被选候选合计只能占用它一次；这一规定本身没有表达该无人机能否先后执行多个航次。

为讨论与巡检模型的关系，首先假定两层使用同一个有限候选集合，并且$x_e=1$在两层中指向同一项执行要求。将补充的执行决策统记为$y$，其内容包括位置轨迹$q$、事件时刻向量$\boldsymbol t$、电量状态$\boldsymbol E$以及任务分配、执行顺序、设备身份和姿态等变量。这里$q$的空间坐标以$\mathrm m$计，$\boldsymbol t$的分量以$\mathrm s$计，$\boldsymbol E$的分量以$\mathrm{Wh}$计；离散身份和分配指示量无量纲。它们是变量块，不能将整个$y$理解为具有同一单位的标量。

设$t_0,T\in\mathbb R$分别为给定班次的起点和终点，满足$t_0<T$，采用同一时钟并以$\mathrm s$计。对每一个需要描述位置的运动对象，其位置轨迹是定义在班次区间上的函数：
\begin{equation}
 q:[t_0,T]\longrightarrow\mathbb R^3,
 \qquad t\longmapsto q(t).
 \label{s0103:trajectory_type}
\end{equation}
式中，$t$是以$\mathrm s$计的时间自变量，$q(t)$是以$\mathrm m$计的三维位置向量。这里用$q$代表单个对象的轨迹，多个对象分别具有各自的轨迹函数；函数所需的连续性或可微性由具体动力学约束规定。轨迹函数与单个事件位置的类型不同：前者描述整个时间区间上的位置变化，后者才是$\mathbb R^3$中的一个向量。

记$\mathcal D$为执行变量各自取值域的乘积，其中包含上述轨迹的函数空间、位于$[t_0,T]$内的事件时刻、非负电量、二元分配指示量以及属于相应有限设备集合的身份变量。姿态变量的表示方式及取值域由运动和相机模型具体规定。道路、禁飞区域、电量上下限等进一步限制由约束施加。用逻辑命题$\Phi(x,y)$表示给定初态、环境信息和参数下，全部物理约束与逻辑约束同时成立，则联合可行域及其在选择变量上的投影定义为
\begin{align}
 \mathcal F_{\mathrm j}
 &=\left\{(x,y)\in\{0,1\}^{n}\times\mathcal D:
                \Phi(x,y)\right\},
 \label{s0103:joint_domain}\\
 \pi_x(\mathcal F_{\mathrm j})
 &=\left\{x\in\{0,1\}^{n}:\exists y\in\mathcal D,
                \ (x,y)\in\mathcal F_{\mathrm j}\right\}.
 \label{s0103:projection}
\end{align}
这里的$\Phi$是对完整约束交集的记号，并非以一个抽象条件替代后续各章的具体建模。对于给定$x$，满足$\Phi(x,y)$的$y$称为该选择的执行见证。它必须同时说明任务如何完成、设备如何连续运动、资源如何周转以及终态如何满足要求。仅列出被选任务，或者分别为各任务找到一条孤立可行路径，均不足以提供共同执行的见证。

\subsection{投影关系成立的条件}

组合可行域与联合可行域的投影是否具有包含关系，取决于资源定义是否忠实保留了执行过程中的必要限制。若已经证明每一个联合可行方案都满足式\eqref{s0103:combinatorial_domain}中的容量约束，即
\begin{equation}
 \forall (x,y)\in\mathcal F_{\mathrm j},\qquad Ax\leq b,
 \label{s0103:necessary_capacity}
\end{equation}
则由投影定义立即得到
\begin{equation}
 \pi_x(\mathcal F_{\mathrm j})\subseteq\mathcal F_{\mathrm c}.
 \label{s0103:outer_domain}
\end{equation}
证明只需取任意$x\in\pi_x(\mathcal F_{\mathrm j})$：它存在一个联合可行的执行见证，因而满足式\eqref{s0103:necessary_capacity}，也就属于$\mathcal F_{\mathrm c}$。在这种情况下，组合模型保留了部分必要条件，可以作为选择层的放松；通过组合检查的选择仍可能无法安排会合、无法获得合格图像，或者没有足够能量返航。因此，组合可行性不自动推出巡检可执行性。

反方向的包含关系需要另一个独立条件：组合模型接受的每个选择都能补全为联合可行方案，即
\begin{equation}
 \forall x\in\mathcal F_{\mathrm c},\quad
 \exists y\in\mathcal D\ \text{使得}\ \Phi(x,y).
 \label{s0103:witness_condition}
\end{equation}
这一条件保证$\mathcal F_{\mathrm c}\subseteq\pi_x(\mathcal F_{\mathrm j})$。若式\eqref{s0103:necessary_capacity}和式\eqref{s0103:witness_condition}同时成立，才能得出
\begin{equation}
 \pi_x(\mathcal F_{\mathrm j})=\mathcal F_{\mathrm c}.
 \label{s0103:conditional_equality}
\end{equation}
这里的等式是有明确前提的集合关系，不是本书现有组合实验已经证明的工程模型等价性。逐个候选拥有可行航迹，也不能替代式\eqref{s0103:witness_condition}，因为多个候选的见证可能要求同一车辆在同一时刻位于不同地点，或者重复使用同一组尚未完成充电的物理电池。

资源定义过于保守时，式\eqref{s0103:necessary_capacity}本身也可能不成立。考虑两个不同的候选$e,f\in E$，二者都需使用同一架无人机。若把该无人机定义为整个班次只能使用一次的容量一资源，组合模型要求
\begin{equation}
 x_e+x_f\leq1.
 \label{s0103:whole_shift_conflict}
\end{equation}
但执行模型可能允许该无人机先完成$e$，经过必要的转移和服务后再完成$f$。若$s_f,h_e\in[t_0,T]$分别为$f$的开始时刻和$e$的结束时刻，$d_{ef}\in\mathbb R_{\geq0}$为两者之间所需的转移与服务时长，三者均以$\mathrm s$计，则这种先后执行至少需要满足
\begin{equation}
 s_f\geq h_e+d_{ef}.
 \label{s0103:sequential_reuse}
\end{equation}
式\eqref{s0103:sequential_reuse}仅说明时间复用所需的一项条件；还须验证位置接续、电池和其他约束。若全部条件确实可满足，$x_e=x_f=1$便可能具有执行见证，却被式\eqref{s0103:whole_shift_conflict}排除。该符号示意说明：把可复用设备当作一次性资源，所得组合模型可能限制了原问题，不能未经证明便称为原问题的放松。

若候选预先固定了占用时间区间，可以考虑将设备与时间区间的组合作为资源；此时还须保证时间划分没有遗漏重叠占用，且候选之间的转移和边界接续得到检查。若占用时刻随路径和调度变化，关联关系就不能在没有额外处理的情况下作为固定矩阵$A$输入。由此，建立两层联系的第一步是统一候选语义和资源语义，而不是仅让两层拥有同名的二元变量。

\subsection{候选内容与评价目标的变化}

纯组合模型给定固定正权重$w_e\in\mathbb R_{>0}$，以无量纲收益分值评价候选，最大化$\sum_{e\in E}w_ex_e$。候选是否值得选择由固定收益与资源竞争共同决定。巡检任务的执行方式则与轨迹、时序和设备状态相联系：同一观测要求可以由不同无人机、不同观测位置和不同航次完成，其时间与用电代价并不必然相同。即使先生成一批可行候选，它们也只是原联合决策空间中已经列举的执行方式。

表\ref{s0103:layer_comparison}列出从组合层进入巡检层需要补充的主要信息。表中组合层一列描述基本固定关联模型的表达范围，并不排除通过扩充候选与资源定义建立更复杂的有限模型。

\begin{table}[htbp]
 \centering
 \small
 \caption{组合选择与巡检执行之间的信息联系}
 \label{s0103:layer_comparison}
 \begin{tabularx}{\linewidth}{@{}p{0.12\linewidth}XXX@{}}
  \toprule
  对象 & 组合层 & 巡检层 & 需要新增的信息\\
  \midrule
  候选任务 & 内容、资源需求和收益预先固定 & 观测要求与设备分配、航次及路径共同确定执行方式 & 目标位置、任务完成规则与允许的执行方式\\
  时间复用 & 以固定容量累计占用份数 & 同一设备可以在满足接续条件后再次使用 & 占用起止时刻、服务时长与前后状态\\
  车辆会合 & 资源关联不直接给出会合过程 & 车辆与无人机须在兼容的地点和时刻完成起降或回收 & 道路可达性、车辆运动、停靠位置与操作窗口\\
  成像观测 & 选中候选表示接受既定任务 & 完成任务还须满足可见性和图像质量要求 & 目标几何、相机参数、观测距离、姿态与遮挡信息\\
  飞行能量 & 一份资源占用不表示实际耗电 & 用电随运动、持续时间和环境条件变化 & 功率关系、天气、初始电量与安全余量\\
  电池与充电 & 固定关联不能单独描述电量恢复过程 & 物理电池按身份周转，充电占用时间和通道 & 电池身份、容量、充电功率、效率和可用通道\\
  \bottomrule
 \end{tabularx}
\end{table}

评价目标也需要相应区分。令$\mathcal G$为有限杆塔集合，$z_g\in\{0,1\}$表示杆塔$g\in\mathcal G$的全部规定观测是否完成并满足验收要求。对一个给定信息情形，巡检模型依次评价未完成杆塔数、全队返库经过时间和机载用电：
\begin{equation}
 \operatorname{lexmin}\left(
      \sum_{g\in\mathcal G}(1-z_g),\ C-t_0,\ E_{\mathrm{air}}
                         \right).
 \label{s0103:inspection_objective}
\end{equation}
其中，$C\in[t_0,T]$为全队返库完成时刻，包含车辆及其回收的无人机，采用与$t_0$相同的时钟并以$\mathrm s$计；$E_{\mathrm{air}}\in\mathbb R_{\geq0}$为纳入当前优化范围的机载用电，单位为$\mathrm{Wh}$。第一项是无量纲整数，取值为$0$至$|\mathcal G|$。词典序要求先最小化第一项，再在第一项最优的方案中最小化第二项，最后在前两项均最优的方案中最小化第三项；它不把不同单位的三项直接相加，也不将联合决策拆成三个互不关联的物理问题。机载用电不包含车辆牵引能耗，充电电网输入另行报告。若以$\mathrm W$为单位的功率对以$\mathrm s$为单位的时间积分，所得数值须除以$3600$才能换算为$\mathrm{Wh}$；滚动重规划时，已经发生的用电是常数，不重复计入待优化的用电。

这一目标允许部分杆塔未完成，而安全、成像质量、能源下限和会合条件仍须满足。某一固定候选集合无法实现全部杆塔覆盖，不能据此断言允许自由选择轨迹和时序的原联合问题无解；也不能把原目标改写成所有杆塔必须完成的硬约束。反过来，即使两个模型具有相同的选择投影，只要评价函数及其优先级不同，最优选择也未必相同。

\subsection{超图关联表示的作用与证据边界}

两层模型均可使用超图描述关联，但顶点与超边的含义必须分别说明。在式\eqref{s0103:combinatorial_domain}中，顶点是资源，超边是候选任务；选择超边就是接受任务。在联合模型的有限数值转录中，还可以采用另一种表示：将标量决策变量作为顶点，将每条约束实际依赖的变量集合表示为超边。

具体地，令$\mathcal W$为有限转录中的标量决策变量索引集合，$\mathcal C$为约束索引集合。对于每个$c\in\mathcal C$，定义其作用域$S_c\subseteq\mathcal W$及相应的约束依赖超图：
\begin{align}
 S_c&=\left\{i\in\mathcal W:
       \text{第$c$条约束实际依赖第$i$个决策变量}\right\},
       \qquad c\in\mathcal C,
 \label{s0103:constraint_scope}\\
 H_{\mathrm{dep}}
 &=\left(\mathcal W,(S_c)_{c\in\mathcal C}\right).
 \label{s0103:dependency_hypergraph}
\end{align}
$\mathcal W$、$\mathcal C$及$S_c$都是无量纲索引集合，变量本身仍保留各自的取值域和物理单位。这里采用带约束索引的超边族$(S_c)_{c\in\mathcal C}$：即使两个不同约束具有相同的作用域，也保留各自的约束身份，不将它们合并为一条约束。所谓实际依赖，是指该变量参与决定约束的取值或逻辑真假，不是指它在某个计算点的导数恰好非零。目标函数的依赖另行登记。

例如，一条会合约束可以同时依赖车辆位置、无人机位置和会合时刻；一条电池接续约束可以同时依赖电池身份、前后航次、电量与充电时长。式\eqref{s0103:dependency_hypergraph}记录这些变量共同作用于哪条约束，却没有保存足以替代约束方程的全部信息。同一依赖集合可以对应不同的函数、不同的上下界和不同的可行域。多个航次读取同一项外生天气数据，也不表示它们竞争一份可耗尽的资源。

因此，约束依赖超图不能直接当作资源匹配超图使用。前者的超边表示必须满足的关系，后者的超边表示可以接受或拒绝的候选；二者承担的数学作用不同。辅助变量的引入、变量消元及标量或变量块粒度的选择还会改变依赖超边的大小与数量。比较这些结构量时，必须先固定转录方式和计数口径，不能用观测任务数代替标量决策变量数。

当前全部纯组合实验采用容量一，其二元可行域能够由两两冲突表达。较大的超边秩只说明一个候选涉及较多资源，不能单独证明不可约的高阶耦合。对于巡检联合模型，更细的依赖表示有助于追踪约束连接了哪些决策，但不会自动给出可用的求解算法，也不会自动产生速度或质量优势。性能结论仍需在相同物理约束、信息截止、初末状态和目标下进行实际比较。纯组合实验能够检验固定候选及固定关联下的选择方法；将所得选择用于巡检，还须补足并核验相应的执行见证。
